% 使用 XeLaTeX 编译两次。
\documentclass[a4paper,11pt]{article}
\usepackage[UTF8]{ctex}
\usepackage[margin=2.1cm,headheight=15pt]{geometry}
\usepackage{booktabs,tabularx,array,enumitem,float,amsmath}
\usepackage{xcolor,tikz,fancyhdr,hyperref}
\usetikzlibrary{arrows.meta,positioning,calc}
\definecolor{navy}{HTML}{244C67}
\definecolor{pale}{HTML}{F0F5F8}
\hypersetup{colorlinks=true,linkcolor=navy,urlcolor=navy,
  pdftitle={挑战赛道技术路线与五人任务分工},pdfauthor={项目组}}
\newcolumntype{L}{>{\raggedright\arraybackslash}X}
\renewcommand{\arraystretch}{1.2}
\linespread{1.08}
\setlist[itemize]{leftmargin=1.6em,itemsep=3pt,topsep=4pt,parsep=0pt}
\setlist[enumerate]{leftmargin=1.8em,itemsep=4pt,topsep=4pt,parsep=0pt}
\setlength{\parskip}{3pt}
\setlength{\emergencystretch}{2em}
\pagestyle{fancy}
\fancyhf{}
\fancyhead[L]{\small 挑战赛道技术路线与任务分工}
\fancyhead[R]{\small XH-202602}
\fancyfoot[C]{\thepage}
\newcommand{\route}[1]{\noindent\textbf{参考路线：}#1\par}

\begin{document}
\title{0908任务分配\\[0.4em]\large ——挑战赛道技术路线与任务分工}
\author{}
\date{}
\maketitle
\vspace{-1.8em}
\section{挑战赛道要求概括}
\subsection{挑战目标与指定平台}
挑战赛道在基础赛道能力之上，重点验证\textbf{模型轻量化、指定芯片适配、异构计算优化及性能保持}。根据晋级通知，统一采用\textbf{地平线J6P}，支持先在PC端使用X86仿真调试，最终以Docker镜像提交。

本组沿用已经形成的模块化VLA链路，不重新训练大型端到端VLA。继续使用现有RTX 5090服务器完成训练与CARLA联调，使用独立的J6P官方工具链完成量化、编译和仿真；板端时延、功耗、内存及利用率由J6P实机验证。以下按任务依赖组织，不安排具体日程。

\subsection{核心要求与评分目标}
\noindent\begin{tabularx}{\textwidth}{p{2.3cm}p{4.0cm}L}
\toprule
维度 & 题目要求或高分目标 & 落实与统计方式\\
\midrule
性能保持 & 核心性能下降不超过3\% & 压缩前后采用同输入、同任务与同统计方法；继续满足基础场景最低要求。\\
推理响应 & 模型端到端推理不超过200\,ms & 保留基础场景更严格的响应要求；逐请求测量，不只测单个网络。\\
计算量压缩 & 压缩后/原模型FLOPs不超过0.5 & 通过结构与特征压缩降低计算量，量化位宽变化不直接视为FLOPs减半。\\
资源占用 & 功耗不超过25\,W，内存不超过512\,MB & 统计实际运行功耗和峰值占用，不能用权重大小代替运行内存。\\
异构执行 & 对应评分档位利用率目标不低于80\% & 依据真实有效负载测量调度效果，不增加无效运算来提高利用率。\\
验证与材料 & 同源评测、部署模型与Docker & 准备三场景1000帧代表性多模态数据及对照结果，保留优化和部署证据。\\
\bottomrule
\end{tabularx}

\clearpage
\section{挑战赛道结构划分与人员职责}
\subsection{总体结构}
\begin{figure}[H]
\centering
\begin{tikzpicture}[
box/.style={draw=navy,fill=pale,align=center,text width=6.5cm,minimum height=1.05cm,font=\small,inner sep=5pt},
wide/.style={box,text width=13.6cm},
arr/.style={-{Latex[length=2.2mm]},thick,draw=navy}]
\node[box] (asr) at (-3.7,0) {语音识别与文本规范化\\李畅锦：ASR适配、驾驶语义保持};
\node[box] (sensor) at (3.7,0) {多视角图像、LiDAR、车辆状态\\刘旭：采集、同步与回放};
\node[box] (lang) at (-3.7,-1.8) {轻量语言编码与结构化指令解析\\朱善哲：DrivingIntent + 共享Embedding};
\node[box] (vision) at (3.7,-1.8) {视觉/BEV编码与候选实体\\黄皓星：感知压缩、场景理解与语义对齐};
\node[wide] (adapter) at (0,-3.65) {轻量Cross-Attention Decision Adapter + 历史状态\\朱善哲主责；黄皓星提供场景/实体特征；王皓然约定动作语义};
\node[wide] (gate) at (0,-5.4) {风险安全门、条件执行与有限状态机（FSM）\\王皓然：动作仲裁、目标速度、时序一致性与安全兜底};
\node[wide] (control) at (0,-7.15) {CARLA控制、执行反馈与统一评测\\刘旭：闭环运行、日志、资源测量、指标与录屏};
\draw[arr] (asr)--(lang);
\draw[arr] (sensor)--(vision);
\draw[arr] (lang.south)--(adapter.north -| lang.south);
\draw[arr] (vision.south)--(adapter.north -| vision.south);
\draw[arr] (adapter)--(gate);
\draw[arr] (gate)--(control);
\end{tikzpicture}
\caption{保留基础赛道链路，在原职责边界内开展轻量化与J6P适配。}
\end{figure}

\subsection{人员职责总览}
\noindent\begin{tabularx}{\textwidth}{p{1.6cm}p{3.5cm}L}
\toprule
负责人 & 延续的工作方向 & 挑战赛道职责\\
\midrule
朱善哲 & 指令解析、轻量VLA、集成 & 语言与Adapter压缩、共享编码、公共J6P后端、版本整合和整体性能优化。\\
李畅锦 & 语音识别与预处理 & ASR轻量适配，驾驶语义、数字方向词、方言噪声和音频链路验证。\\
黄皓星 & 场景理解与语义对齐 & 视觉/BEV压缩，候选实体与目标落地，感知侧导出、量化和精度验证。\\
王皓然 & 高层动作、风险门与FSM & 动作接口、时序安全、条件执行、异常降级和规则侧运行优化。\\
刘旭 & CARLA、控制、日志与评测 & 同源数据、传感器回放、控制闭环、统一指标、板端资源测量与重复验证。\\
\bottomrule
\end{tabularx}

\clearpage
\section{各负责人需要完成的任务}
\subsection{朱善哲：指令解析、VLA优化与系统集成}
\begin{itemize}
\item 优化ModernBERT与轻量VLA，减少重复编码和计算开销，保持指令解析与决策能力。
\item 完成语言和决策模型的J6P适配，协调各模块接口与整体联调。
\end{itemize}
\route{共享语言编码，结合蒸馏、结构压缩和量化进行优化，再比较精度与全链路性能。}

\subsection{李畅锦：语音识别优化}
\begin{itemize}
\item 降低语音识别模块的资源占用与延时，保持驾驶指令识别效果。
\item 检查方言、噪声以及方向、数字等关键信息，并与指令解析模块联调。
\end{itemize}
\route{先测试现有模型开销，再尝试量化或轻量模型替换，以驾驶语义保持为主要优化依据。}

\subsection{黄皓星：场景理解与语义对齐优化}
\begin{itemize}
\item 优化视觉与BEV编码，降低感知计算量，保留关键目标和场景信息。
\item 保持指令目标与场景实体的准确对齐，完成感知侧J6P适配。
\end{itemize}
\route{先量化现有模型，再按效果进行剪枝、蒸馏或特征缩减，结合下游决策验证优化收益。}

\subsection{王皓然：高层决策与安全控制优化}
\begin{itemize}
\item 完善高层动作、风险安全门和FSM，保证组合指令与条件指令正确执行。
\item 保持连续决策稳定，处理危险建议、推理超时和感知异常等情况。
\end{itemize}
\route{延续现有决策架构，重点回归加减速波动、前车制动等问题，再优化状态管理和安全策略。}

\subsection{刘旭：仿真控制与统一评测}
\begin{itemize}
\item 维护三类CARLA场景与控制链路，组织原版和优化版的重复对照测试。
\item 配合J6P部署，测量任务效果、稳定性、延时和资源占用，定位系统瓶颈。
\end{itemize}
\route{使用同源数据回放和多次闭环测试比较效果，再通过J6P实机验证资源与运行性能。}
\end{document}
